\ficha{Almeida (2010), Política econômica e crises cambiais no Império}{almeida2010}

\begin{identificacao}
ALMEIDA, José Tadeu de. Política econômica e crises cambiais: a gestão
financeira do Império do Brasil nos primórdios do padrão-ouro (1846-1858).
\textit{História e Economia Revista Interdisciplinar}, v. 7, n. 1, p. 49-69,
2010.

Artigo de revista, em português, lido na íntegra: p. 49 a 66 de texto com sete
gráficos e p. 67 a 69 de bibliografia. O autor assina como doutorando em
História Econômica na FFLCH-USP. O texto deriva da dissertação de mestrado dele
na Unicamp (2010) e se autodenomina ``dissertação'' na conclusão (p. 66).
\end{identificacao}

\bloco{Problema e tese}

Como uma economia periférica administra a adesão ao padrão ouro-libra sem banco
central, com sistema bancário incipiente e dependência do crédito inglês?
Almeida responde pela gestão conservadora dos gabinetes da Conciliação, entre
1846 e 1858. O foco no equilíbrio orçamentário, na valorização cambial e na
recusa da emissão não vem, para ele, de doutrina metalista, e sim do interesse
estratégico de construir o Estado nacional sob a restrição externa de conservar
legitimidade para a paridade adotada.

\bloco{Argumento}

\begin{enumerate}[leftmargin=*,nosep]
\item A Lei nº 401, de 11 de setembro de 1846, fixa a paridade em 27 pence por
mil-réis e, no artigo segundo, autoriza o Governo Imperial a recolher
papel-moeda na quantidade que julgar necessária para manter a valorização da
moeda (p. 55). O padrão-ouro entra no Brasil junto com um instrumento de
contração, não apenas como regra de conversão.

\item A ortodoxia vem de fora, do interesse dos credores.

\chave{tal sistema monetário, gerido no âmbito internacional, visava, sobretudo,
salvaguardar interesses dos grandes credores (notoriamente as casas bancárias
inglesas)}{p. 52}

\item A hipótese central é uma contradição, não um alinhamento: a expansão dos
negócios internos pede crédito e emissão, e a legitimidade internacional da
paridade pede defesa do mil-réis e restrição à emissão livre (p. 51). Os
ministros da Fazenda acomodam os dois lados em vez de escolher um.

\item A crise de 1857 mede o custo da inserção periférica. Pressionado por
credores a liquidar débitos, o governo suspende a conversibilidade entre notas
do Tesouro e moeda metálica em 11 de novembro (p. 61). A saída não é o câmbio,
é o crédito externo: o Império avaliza empréstimos ao Banco do Brasil de
\textsterling{}340.000 e depois de \textsterling{}600.000 junto à Rothschild
\& Sons, descrita como emprestador de última instância (p. 62). O mil-réis cai
a 23,5 pence em dezembro e volta a 26,55 pence em agosto de 1858 (p. 63).

\item A inflexão de 1858-59 restaura a ortodoxia por meio de pessoas, não de
instituições. Sousa Franco, liberal na Fazenda em 1857, priorizou o meio
circulante e a pluralidade emissora. Torres Homem, seu sucessor, ``em um corte
nitidamente metalista'', atribui ao papel-moeda a alta de preços, a
desvalorização cambial e a queda dos salários (p. 63), e suspende a autorização
de 1856 que permitia ao Banco do Brasil circular notas no triplo das reservas
(p. 64).

\item A conclusão recusa a taxonomia de época como chave explicativa.

\chave{por trás das posturas aparentemente discricionárias, porque
`emissionistas', ou draconianas, porque restritivas, realmente permeava a ação
daqueles estadistas (...) a da construção e afirmação do novel Estado
Nacional}{p. 66}
\end{enumerate}

\bloco{Conceitos e definições operacionais}

\textbf{Padrão-ouro}: instituição de um momento histórico, na leitura tomada de
Eichengreen, em que toda prioridade governamental se vincula a manter reservas
compatíveis com a paridade exigida (p. 50).

\textbf{Inserção periférica}: a periferia é rotineiramente obrigada a rever ou
suspender suas metas de paridade-ouro, enquanto o sistema se mantém em poucos
países-chave (p. 56). O Brasil é conversível de jure e vulnerável de facto (p.
65).

\textbf{Papel-moeda e moeda-papel}: papel-moeda é emissão com lastro inferior a
100\%; moeda-papel é certificado com lastro integral, e só este é compatível
com os preceitos do padrão-ouro (p. 55, nota 10).

\textbf{Metalista e papelista}: rótulos derivados da postura mais ou menos
restritiva quanto à liberdade de emissão, insuficientes, para o autor, para
explicar a ação dos ministros (p. 66).

\bloco{Evidência e método}

Narrativa de política econômica sobre bibliografia secundária e fontes citadas
de segunda mão, sobretudo Carvalho (1927), Mont'Alegre (1972), Calógeras (1960)
e Villela (1999). Sete gráficos cobrem comércio, meio circulante, reservas e
câmbio. Não há arquivo primário nem estimação. Comparação pontual: em 1857 os
bancos franceses podiam emprestar até 100\% das reservas, e no Brasil a razão
chegou a 300\% (p. 57).

\bloco{Agentes e vertentes}

No campo restritivo, Itaboraí, Carneiro Leão e Salles Torres Homem, do Partido
Conservador. No campo emissionista, Bernardo de Sousa Franco, liberal, e Mauá.
Do lado credor, Rothschild \& Sons e o Banco da Inglaterra. Almeida se apoia em
Eichengreen para a definição do padrão, em Kindleberger para a mecânica do
pânico de 1857, e em Gremaud e Villela para a história monetária brasileira.

\bloco{Críticas e limites}

Densidade probatória baixa e aparato com falhas: Eichengreen aparece com dois
anos para o mesmo livro, e Villela é grafado de três maneiras. O autor admite
que tratar o padrão-ouro como indutor de investimento no Brasil oitocentista
seria ``hipótese heróica'' sem indícios suficientes (p. 65, nota 25). A tese da
construção estatal é afirmada, não demonstrada contra a leitura doutrinária
concorrente. O recorte termina em 1858 e nada alcança 1898-1906, a Caixa de
Conversão ou a imprensa.

\begin{dialogo}
\item Procurar, em 1906, a invocação do par legal de 27 pence e da Lei de 1846
como padrão de verdade monetária. Vocabulário: ``par legal'', ``27 pence'',
``lei de 1846'', ``reforma monetária do Império''. Testa se a ortodoxia de 1906
se legitimava por precedente imperial ou pelo funding loan de 1898.

\item Rastrear os nomes de Sousa Franco, Itaboraí, Torres Homem e Mauá no
debate de 1906-1914. Testa se os agentes construíam genealogia imperial para
suas posições.

\item Ler a cobertura da suspensão do troco em 1914 procurando referência à
suspensão da conversibilidade de novembro de 1857. Testa se havia memória
pública de que suspender era resposta periférica recorrente, e não ruptura
inédita.
\end{dialogo}

\usonocapitulo{Parte III, como apoio, para uma afirmação precisa: a ortodoxia
de 1898-1906 tem precedente imperial verificável, porque a paridade de 27 pence
de 1846 já vinha com autorização legal para recolher papel-moeda (p. 55), e a
inflexão de 1858-59 já reunia equilíbrio orçamentário, valorização cambial e
recusa da emissão em um mesmo programa (p. 63-64). Serve também para
relativizar a taxonomia de época, já que o texto sustenta que metalista e
papelista são rótulos derivados da postura quanto à emissão e não explicam a
ação dos ministros (p. 66), o que apoia o eixo ortodoxo contra expansionista
adotado pela dissertação. Como o recorte se encerra em 1858, o artigo não
sustenta afirmação alguma sobre a Caixa de Conversão e entra apenas como
antecedente.}
